\section*{Chapter \ref{ch:iv2:portfolio-construction} --- Portfolio Construction}
\markboth{\protect\omTocMark{}{Answers: Portfolio Construction}}{\protect\omTocMark{}{Answers: Portfolio Construction}}

\begin{solution}{iq:iv2:portfolio-construction:1}
The weights are about 85.7\% and 14.3\%, and the volatility is about 19.6\%.

$w = (\sigma_2^2 - \sigma_{12})/(\sigma_1^2 + \sigma_2^2 - 2\sigma_{12})$, with $\sigma_{12} = 0.03$. So $w = 0.06/0.07 = 6/7$.

The variance is $w^2(0.04) + (1 - w)^2(0.09) + 2w(1 - w)(0.03) = 0.0386$.

\iqwrong{Allocating by inverse variance alone. That ignores the correlation and gives 69/31, which has a higher variance here.}
\end{solution}

\begin{solution}{iq:iv2:portfolio-construction:2}
It is 50/50, with a volatility of about 15.5\%.

By symmetry, the weights are equal, and the variance is $\sigma^2(1 + \rho)/2 = 0.04 \times 0.6 = 0.024$. Its root is 15.49\%.

\iqwrong{Saying that correlation changes the weights. With equal volatilities, it only changes the resulting risk.}
\end{solution}

\begin{solution}{iq:iv2:portfolio-construction:3}
The weights are 2/3 and 1/3, and no, the correlation doesn't matter for two assets.

A contribution is $w_i(\Sigma w)_i$. For two assets, equating them gives $w_1^2\sigma_1^2 + \sigma_{12}w_1w_2 = w_2^2\sigma_2^2 + \sigma_{12}w_1w_2$. The cross term cancels, so $w_1\sigma_1 = w_2\sigma_2$.

\iqwrong{Carrying that over to three or more assets. There, inverse volatility is equal risk only if all correlations are equal.}
\end{solution}

\begin{solution}{iq:iv2:portfolio-construction:4}
It is 0.50 with 100 bets, and about 0.79 with 252.

$\mathrm{IR} = \mathrm{IC}\sqrt{\mathrm{BR}}$: $0.05 \times 10 = 0.5$, and $0.05 \times 15.87 = 0.79$.

\iqwrong{Counting trades as bets. Breadth counts independent forecasts per year, and a stock held across 252 days on one slow signal is not 252 bets.}
\end{solution}

\begin{solution}{iq:iv2:portfolio-construction:5}
One-way turnover is 10\%, and it costs 2 basis points, or about 24 basis points a year.

The sum of absolute changes is 0.2, and half of that is the one-way 0.1. You trade 0.2 of a dollar, sell 0.1 and buy 0.1, at 10 basis points: $0.2 \times 10 = 2$ basis points of portfolio. Twelve times gives 24.

\iqwrong{Costing the one-way number. You pay on both the sale and the purchase, and quotes of turnover vary, so check which is meant.}
\end{solution}

\begin{solution}{iq:iv2:portfolio-construction:6}
You short \$15 million, and your net dollar position is $-\$5$ million.

The long has beta dollars of $10 \times 1.2 = 12$. Then $12/0.8 = 15$.

\iqwrong{Assuming that beta neutral and dollar neutral are the same. Here, being neutral on one means being \$5 million short on the other. If the beta estimate is off, so is the hedge.}
\end{solution}

\begin{solution}{iq:iv2:portfolio-construction:7}
The tangency weights are about 72.4\% and 27.6\%, and the Sharpe is about 0.429. At a risk aversion of 4, they hold 38.9\% and 14.8\%, with a return of about 4.59\% and a volatility of about 10.7\%.

Solve $\Sigma x = \mu$. Here $x = (1.556, 0.593)$. Normalising gives the tangency weights. The Sharpe is $\sqrt{\mu^\top x} = \sqrt{0.1837} = 0.429$.

The investor holds $x/\lambda$, so $(0.389, 0.148)$, and the rest is in cash. So the return is $0.1837/4 = 4.59\%$, and the volatility is $0.429/4 = 10.7\%$.

\iqwrong{Forcing weights to sum to 1 for the investor. With a risk-free asset, risk aversion sets the size: it scales the whole risky portfolio, and the rest is cash.}
\end{solution}

\begin{solution}{iq:iv2:portfolio-construction:8}
The weights are about 123.5\% and $-23.5\%$, with a volatility of about 19.5\%. The second weight is negative when $\rho > \sigma_1/\sigma_2$, and here $0.8 > 0.667$.

The covariance is $0.048$, so $w = (0.09 - 0.048)/(0.13 - 0.096) = 1.235$.

The second weight is $(\sigma_1^2 - \rho\sigma_1\sigma_2)$ over a positive denominator. It's negative when $\sigma_1 < \rho\sigma_2$. At high correlation, the riskier asset is mostly a levered copy of the other. So shorting it hedges the first.

\iqwrong{Treating a short in a minimum-variance portfolio as an error. It's the correct hedge, but it does depend on an estimate of correlation that may not hold in a stress.}
\end{solution}

\begin{solution}{iq:iv2:portfolio-construction:9}
The Sharpe is about 0.674, with $x = (5.91, -2.73)$. The best long-only Sharpe is 0.50, all in A.

Solve $\Sigma x = \mu$, with variances 0.04 and 0.09 and covariance 0.05. So $x = (5.91, -2.73)$, and the Sharpe is $\sqrt{\mu^\top x} = 0.674$.

For long-only, the unconstrained solution wants B negative. The Sharpe is unimodal along the mix of A and B, so the long-only best is at the boundary, with no B. And $0.10/0.20 = 0.5$.

So being able to short B would add 0.17 of Sharpe, but at a big leverage: 5.91 long and 2.73 short per unit of risk aversion, on a covariance that has to be right.

\iqwrong{Saying a high-return asset must be in any good portfolio. B has a positive return and is still shorted here, because it's highly correlated with A and has a worse Sharpe.}
\end{solution}

\begin{solution}{iq:iv2:portfolio-construction:10}
It is 18.2\%, 17.1\% and 16.6\%, and the limit is 16.4\%. You need about 19 stocks to be within 1 point.

The variance is $\sigma^2(1/n + (1 - 1/n)\rho)$. As $n \to \infty$, $\sigma\sqrt{\rho} = 30\% \times 0.548 = 16.43\%$.

Within 1 point means below 17.43\%. So $0.09(0.7/n + 0.3) \le 0.1743^2 = 0.03038$, and $n \ge 19$.

\iqwrong{Saying more stocks always means much less risk. The first 19 do nearly all the work, and the common factor stays, however many you add.}
\end{solution}

\begin{solution}{iq:iv2:portfolio-construction:11}
Effective breadth is about 83, and the information ratio about 0.18. Treated as independent, it's about 0.45.

With $N$ bets at pairwise correlation $\rho$, the variance of the average is $(1 + (N - 1)\rho)/N$, so the effective number is $N/(1 + (N - 1)\rho) = 500/5.99 = 83.5$. So $\mathrm{IR} = 0.02\sqrt{83.5} = 0.183$.

Even a correlation of 0.01 cuts breadth by a factor of six. A factor bet disguised as 500 names has this problem.

\iqwrong{Counting stocks for breadth. A small common error across many names reduces the number sharply, and at the limit gives $1/\rho$ independent bets, here 100.}
\end{solution}

\begin{solution}{iq:iv2:portfolio-construction:12}
It is about 0.34.

$\mathrm{IR} = \mathrm{TC} \times \mathrm{IC}\sqrt{\mathrm{BR}} = 0.3 \times 0.05 \times 22.36 = 0.335$.

The transfer coefficient is the correlation between the active weights you hold and the weights your forecast alone would give. It measures how much of the signal the constraints let through.

To raise it, relax long-only (a long-only portfolio can only underweight a small stock by its benchmark weight), widen the universe, or reduce the tracking-error limits. A long-short sleeve can raise it most.

\iqwrong{Improving the IC first. Before searching for more signal, check how much you're throwing away in construction: here 70\%.}
\end{solution}

\begin{solution}{iq:iv2:portfolio-construction:13}
The standard error of each mean is about 4.5\%. Shrinking gives 10\% and 0\%. Mean-variance optimisation then gives much less extreme bets, as it maximises the errors less.

$20\%/\sqrt{20} = 4.47\%$. A mean of 15\% has a $t$ of 3.4, and the gap between the two has a standard error of 6.3\%, for a $t$ of 3.2.

Shrinking halves the spread from 20 to 10 points. Weights are linear in the means, and depend on their differences for equal risk. So the long-short tilt halves.

\textbf{Why error-maximising.} The optimiser loads up on assets with high estimated return and low estimated risk. These are likely to be where the estimation error is most in your favour. So in-sample performance overstates, and out-of-sample disappoints.

\iqwrong{Thinking 20 years is a lot of data. The mean depends on the length of time, not the number of observations. High-frequency data helps a lot with variance, but hardly at all with the mean.}
\end{solution}

\begin{solution}{iq:iv2:portfolio-construction:14}
The condition number is 500. After shrinking, the eigenvalues are 4.40, 1.20 and 0.41, and the condition number is about 10.8.

The trace average is $6.01/3 = 2.003$. So each eigenvalue is $0.8\lambda + 0.2 \times 2.003$, with the eigenvectors unchanged.

A minimum-variance portfolio is $\propto \Sigma^{-1}\mathbf{1}$, and the inverse magnifies small eigenvalues. An eigenvalue of 0.01 is usually noise: it says there's a nearly riskless combination, and the optimiser piles into it. Lifting it to 0.41 reduces that weight by a factor of 41.

\iqwrong{Treating a small eigenvalue as real. Its sample estimate is biased low when assets are many relative to data, and with more assets than days, some are exactly zero.}
\iqfollow{How do you choose the shrinkage intensity?}{By minimising expected squared error, which has a closed-form estimate. Or by testing out-of-sample variance of the optimised portfolio.}
\end{solution}

\begin{solution}{iq:iv2:portfolio-construction:15}
With 100 stocks and 1,000 days, the eigenvalues lie in about $[0.47, 1.73]$. With 500, the upper end rises to about 2.91, and the lower to about 0.09.

With $q = N/T$, the edges are $(1 \pm \sqrt{q})^2$. At $q = 0.1$, that's $0.468$ and $1.732$. At $q = 0.5$, it's $0.086$ and $2.914$.

So a noise-only matrix still has an eigenvalue of 1.7 or 2.9, and even a true identity looks structured. Real data has a large market eigenvalue well outside, and perhaps some sectors.

Use it as a filter: eigenvalues within the noise band are treated as noise and replaced by their average, which keeps the trace.

\iqwrong{Using every principal component from a sample as a factor. The ones inside the band carry little beyond noise, and they make a backtest look better than it is.}
\iqfollow{What assumption can this break?}{The bands assume independence through time and thin tails. Fat tails and volatility clustering widen the sample spectrum, so the real cut-off is higher.}
\end{solution}

\begin{solution}{iq:iv2:portfolio-construction:16}
It's 125,250 for the full matrix, and 5,555 for the factor model.

Full: $500 \times 501/2$. Factor: $500 \times 10$ loadings, $10 \times 11/2 = 55$ factor covariances, and 500 idiosyncratic variances.

The assumption is that residuals are uncorrelated once factors are removed. That makes $\Sigma = BFB^\top + D$. It also makes it invertible, even with fewer days than stocks.

It fails where residuals are in fact linked: pair relations, a shared missing factor such as a crowded trade, or a merger. So check residual correlations, and don't treat a low factor risk as low risk when a book is short a common residual.

\iqwrong{Calling factor models an approximation for lack of data only. They also regularise the problem, so they often beat the sample matrix out of sample even with plenty of data.}
\end{solution}

\begin{solution}{iq:iv2:portfolio-construction:17}
Each asset has an expected log growth of $-2\%$ a year, and the rebalanced portfolio about $-1\%$. Buying and holding gets about $-1.2\%$ over 30 years. With a correlation of 1, rebalancing gains nothing.

The log growth is the arithmetic return less half the variance. Each asset has $0 - 0.04/2 = -2\%$. The portfolio has variance $0.02$, so $0 - 0.01 = -1\%$.

So the gain from each individual asset is $\tfrac12(\overline{\sigma_i^2} - \sigma_p^2) = 1\%$, the \index{diversification return}diversification return. It is what you'd get by comparing a balanced portfolio to the average asset.

A buy-and-hold portfolio starts diversified too, but it drifts into the winner, and in the end looks like one asset. So over 30 years, the expected log growth rate is about $-1.2\%$, and the gap to rebalancing is only about 0.2\% a year in this setup. Costs can take that away.

With a correlation of 1, both assets are the same, the portfolio variance is unchanged, and no gain arises.

\iqwrong{Calling rebalancing free money. It raises log growth by lowering the variance drag, not by adding any return: expected final wealth is no higher here, and it is lower when one asset trends.}
\end{solution}

\begin{solution}{iq:iv2:portfolio-construction:18}
The true Sharpe is about 0.57. The estimated tangency shows about 1.65 in sample, and only about 0.19 out of sample. So you've lost two-thirds of the truth while reporting three times the truth.

\textbf{Truth.} For equal weights, return is 0.5\% a month, and volatility is $5\%\sqrt{0.1 + 0.9 \times 0.3} = 3.04\%$. So the monthly Sharpe is 0.164, or 0.569 a year. By symmetry, this is the true optimum.

\textbf{Estimated.} With 10 assets and 60 months, the means' error has the same scale as the means. A sample Sharpe is high because the weights exploit noise. In a simulation of 3,000 runs, the in-sample Sharpe averages 1.65, and the same weights evaluated on the true parameters average 0.19.

So: equal weights or strong shrinkage beat sample optimisation unless the signal is much stronger than this. You need many more observations per asset, or structure on the means.

\iqwrong{Blaming the covariance alone. Most of the damage here comes from the means. If you knew the mean differences were zero, the estimated covariance would still get you much closer to the optimum.}
\end{solution}

\begin{solution}{iq:iv2:portfolio-construction:19}
The implied returns are 9.0\% and 13.5\%. The posterior is 9.5\% and 13.875\%, which gives weights of 65\% and 40\%. As the uncertainty goes to zero, the posterior goes to 10\% and 14.25\%, and the weights to 70\% and 40\%.

\textbf{Prior.} $\Pi = \delta\Sigma w = 2.5\,(0.036, 0.054) = (0.09, 0.135)$.

\textbf{Posterior.} The update is
\[
\mu = \Pi + \tau\Sigma P^\top(P\tau\Sigma P^\top + \Omega)^{-1}(Q - P\Pi).
\]
Here $P = (1, 0)$, the surprise is $0.01$, $\tau\Sigma_{11} = 0.002$, and $\Omega = 0.002$. So asset 1 moves by half the surprise, $+0.005$. Asset 2 moves by $(0.0015/0.004) \times 0.01 = 0.00375$, through correlation.

\textbf{Weights.} Dividing by $\delta\Sigma$ leaves asset 2 at 0.4, and moves only asset 1. That's a known property: only assets in the view move away from market weight. As $\Omega \to 0$, asset 1 goes all the way to its view, and to weight 0.70. As $\Omega \to \infty$, you get the prior and the market portfolio.

\iqwrong{Saying that the second asset's return is unchanged because you had no view on it. Its return moves through correlation, and that's why its weight stays the same.}
\end{solution}

\begin{solution}{iq:iv2:portfolio-construction:20}
Equal weights give shares of 29.8\%, 14.3\% and 56.0\%. Inverse volatility gives 35.2\%, 38.3\% and 26.6\%. The third set gives one third each, so that's the equal-risk portfolio, with a volatility of about 11.6\%.

A risk share is $w_i(\Sigma w)_i/w^\top\Sigma w$. For the last set, $\Sigma w = (0.0162, 0.0087, 0.0211)$ to 4 decimals. So the contributions are all 0.00447, and the sum of the three is the variance, 0.0134.

Inverse volatility ignores correlation. It would be right here only if the correlations were equal. The second asset has a low volatility and is positively correlated with both others, so it gets too much weight on the inverse-volatility rule.

Equal risk does use correlation, and it doesn't use expected returns. That's a real selling point: it's less sensitive to bad inputs than mean-variance. It's still dependent on the covariance, and it tends to be levered to reach a return target.

\iqwrong{Thinking equal risk contributions means equal volatility of holdings. It's about contributions to the total, after correlation.}
\iqfollow{Is there always a solution?}{Yes for a long-only portfolio with positive variance. Minimise $\tfrac12 w^\top\Sigma w - c\sum_i\ln w_i$, and scale to sum to 1: a convex problem with one answer.}
\end{solution}

\begin{solution}{iq:iv2:portfolio-construction:21}
The global minimum-variance portfolio has a return of about 10.2\% and a volatility of about 14.0\%. At a return of 12\%, the weights are about $(0.019, -0.029, 1.010)$, and the volatility is about 15.0\%.

Let $A = \mathbf{1}^\top\Sigma^{-1}\mathbf{1} = 51.29$, $B = \mathbf{1}^\top\Sigma^{-1}\mu = 5.243$, $C = \mu^\top\Sigma^{-1}\mu = 0.6437$, and $D = AC - B^2 = 5.531$.

The minimum is at return $B/A = 10.22\%$ and volatility $1/\sqrt{A} = 13.96\%$.

For target $m$, $w = \Sigma^{-1}(\lambda\mathbf{1} + \gamma\mu)$ with $\lambda = (C - Bm)/D$ and $\gamma = (Am - B)/D$. The variance is $(Am^2 - 2Bm + C)/D$.

\textbf{Two funds.} Every frontier portfolio is a mix of $\Sigma^{-1}\mathbf{1}$ and $\Sigma^{-1}\mu$, with coefficients linear in $m$. So any two frontier portfolios span it, and you don't need to re-solve for each target.

A return of 10\% is below the minimum-variance return, so that point is on the inefficient side: a higher return is available at the same risk.

\iqwrong{Taking a lower target to mean less risk everywhere. Below the minimum-variance return, risk goes back up.}
\end{solution}

\begin{solution}{iq:iv2:portfolio-construction:22}
It is $(0.4, 0.2, 0.4)$, with a volatility of about 14.9\%.

With the third weight capped at 0.4, the others rise, and the first also hits the cap, so the second takes the remainder.

\textbf{Conditions.} With $w^\top\mathbf{1} = 1$ and $w_i \le 0.4$, the conditions say $\Sigma w = \lambda\mathbf{1} - \nu$, with $\nu_i \ge 0$, and $\nu_i = 0$ unless $w_i$ is at the cap. Here $\Sigma w = (0.0232, 0.0300, 0.0174)$ to four decimals. The asset below the cap is the second, so $\lambda = 0.0300$. The others have $\nu_1 = 0.0068$ and $\nu_3 = 0.0126$. Both are positive, so the conditions hold.

Since the problem is convex, those conditions are sufficient: it's the unique best.

Read the multipliers. They're the gain from relaxing each cap. Relaxing asset 3's cap is worth nearly twice relaxing asset 1's.

\iqwrong{Capping the large weight and rescaling the rest pro rata. That gives $(0.517, 0.083, 0.4)$, which breaks the cap on the first asset and doesn't meet the conditions.}
\end{solution}

\begin{solution}{iq:iv2:portfolio-construction:23}
Yes: the total expected return is about 15.5 basis points, against 10 basis points for a round trip. The break-even signal is about 1.3 basis points a day. Held for just one day, you make 2 and pay 10, and that's a loss.

The decay factor is $\rho = 2^{-1/5} = 0.8706$, and the sum of $a\rho^t$ is $a/(1 - \rho) = 2/0.1294 = 15.45$.

The break-even signal has $a/(1 - \rho) = 10$, so $a = 1.29$.

With only 1 day of holding, you'd need a signal of 10 basis points a day.

So holding period matters as much as signal size. It also implies a no-trade zone: with a position on and a small change of signal, you don't trade, since the new information doesn't pay for the cost. A portfolio should move partly towards its target, and the rate depends on the decay and the cost.

\iqwrong{Comparing the 2 basis points with the 5 basis points of cost. The relevant comparison is what the signal will earn in total, net of a round trip. And a slow signal can afford much more trading than its daily number suggests.}
\iqfollow{What if a second signal arrives before the first has decayed?}{Trade the net. Internal crossing of opposite signals saves cost, and that's one of the larger gains from running a combined book.}
\end{solution}

\begin{solution}{iq:iv2:portfolio-construction:24}
The best mix has a Sharpe of about 1.13, and risk weights of about 60\% and 40\%. At a correlation of 0.8, the second strategy adds nothing. At a negative correlation, it adds a lot.

For unit-volatility strategies, solve $Cx = s$, with $C$ the correlation matrix and $s = (1.0, 0.8)$: $x = (0.835, 0.549)$, so risk weights of 60.3\% and 39.7\%. The Sharpe is
\[
\sqrt{s^\top C^{-1}s} = \sqrt{\frac{1 + 0.64 - 2 \times 0.3 \times 0.8}{1 - 0.09}} = 1.129.
\]

The second adds nothing when its weight is zero, which happens when $\rho = S_2/S_1 = 0.8$. Then its returns are explained by the first, at a slope that's already paid for. Note how this isn't at $\rho = 1$.

At $\rho = -0.3$, the Sharpe is $\sqrt{(1.64 + 0.48)/0.91} = 1.53$. So a negative correlation raises it far more than a better Sharpe would.

\iqwrong{Allocating equal capital. Allocation should be by risk, and strategies usually differ a lot in volatility, so equal capital leaves one running the show.}
\end{solution}

\begin{solution}{iq:iv2:portfolio-construction:25}
The results, with weights in order, are in the table. At 10\% volatility, the beta-neutral book has weights $(0.226, 0.122, -0.369)$ and returns 3.20\%. With costs from flat, it trades $(0.142, 0.069, -0.225)$.

\begin{center}
\begin{tabular}{lccccc}
\toprule
Book & Weights & Beta & Return & Volatility & Sharpe \\
\midrule
Free & $(0.244, 0.155, -0.212)$ & 0.205 & 2.77\% & 7.45\% & 0.372 \\
Beta neutral & $(0.145, 0.078, -0.237)$ & 0 & 2.05\% & 6.40\% & 0.320 \\
Both neutral & $(0.119, 0.119, -0.237)$ & 0 & 2.02\% & 6.35\% & 0.318 \\
\bottomrule
\end{tabular}
\end{center}

\textbf{Unconstrained.} $w = \Sigma^{-1}\alpha/\lambda$. The correlation is common, so $\Sigma$ is easy to build, and the solve gives the weights. The return is $\alpha^\top w$, and the beta is $\beta^\top w = 0.205$. So that portfolio has an unintended market bet.

\textbf{Beta neutral.} Replace $\alpha$ by $\alpha - g\beta$, with $g = \beta^\top\Sigma^{-1}\alpha/\beta^\top\Sigma^{-1}\beta = 3.53\%$. That's the alpha less what's explained by beta. Then $w = \Sigma^{-1}(\alpha - g\beta)/\lambda$, and $\beta^\top w = 0$.

\textbf{Dollar neutral too.} Two constraints, two multipliers. With $G$ the matrix of $\beta$ and $\mathbf{1}$, solve $(G\Sigma^{-1}G^\top)a = G\Sigma^{-1}\alpha$. Then $w = \Sigma^{-1}(\alpha - G^\top a)/\lambda$. The Sharpe hardly moves here from 0.320 to 0.318, so this constraint is cheap. Going to beta neutral costs more, from 0.372 to 0.320.

\textbf{Scale.} Multiply by $10\%/6.40\% = 1.562$. Sharpe stays the same.

\textbf{Costs.} With all three positions kept and their signs known, an absolute cost just lowers each $|\alpha_i|$ by $c$: $\alpha' = (0.058, 0.028, -0.038)$. So use the same beta-neutral formula with $\alpha'$, and you get $(0.142, 0.069, -0.225)$. Then check that no position's sign has changed, and every one is still worth trading. The net objective is about 0.94\%. If a position were small enough for its alpha to fall within $\pm c$, you'd stay flat in it, which is the no-trade zone.

\textbf{Checks without a solver.} The first-order conditions: $\alpha' - \lambda\Sigma w - g\beta = 0$ for some $g$, for each position. Beta zero: $1.2 \times 0.142 + 0.8 \times 0.069 - 1.0 \times 0.225 = 0$, to rounding. And the limits: as $\lambda$ rises, everything shrinks in proportion. As costs rise, the smallest position leaves first.

What I'd say about the result: the expected return is only as good as the alpha, so I'd run this with shrunk alphas and show how the book changes. A beta-neutral book with these inputs has an expected Sharpe around 0.3, which is what I'd take from three names. An estimate is what matters, not the digit.

\iqwrong{Solving the unconstrained problem, then hedging out beta with an index. That's another valid portfolio, but it isn't the optimum unless the index's residual has the right properties. The joint solve uses the names themselves to neutralise beta.}
\end{solution}
